% Local editorial additions to VI/28 (Normalization).
% These examples are independent of canonical CP-VI statements and reader bridges.
\par\medskip
\noindent\textbf{Worked examples: integral closure, normalization and local geometry.}
For a domain $A$ with function field $K$, its normalization is the
integral closure $\overline A\subset K$.  An element $b\in K$ is
integral over $A$ if it satisfies a monic polynomial with coefficients
in $A$.  To identify a proposed normalization $B$, one must check
$A\subseteq B\subseteq K$, integrality of $B/A$, and that $B$ is
integrally closed in $K$.  Reduced rings with several components
require their total ring of fractions instead of a single field.

\par\medskip
\noindent\textbf{Worked example 1 (normalizing the cusp).}
Let $A=k[x,y]/(y^2-x^3)\cong k[t^2,t^3]$, with
$x=t^2$ and $y=t^3$.  Since $t=y/x$ lies in $K=\operatorname{Frac}(A)$,
we may compare $A$ with $B=k[t]\subset K$.
The element $t$ satisfies the monic equation
\[
 T^2-x=0,
\]
so $t$ is integral over $A$ and $B=A[t]$ is integral over $A$.
Indeed every polynomial in $t$ reduces using $t^2=x$ to
$a+bt$ with $a,b\in A$, so $B=A+At$ is a finite $A$-module.
The polynomial ring $k[t]$ is a principal ideal domain and hence
integrally closed.  Since it has fraction field $k(t)=K$, it is the
entire integral closure of $A$.
Geometrically the normalization map is
\[
 \mathbb A^1_k\longrightarrow \operatorname{Spec}A,
 \qquad t\longmapsto(t^2,t^3).
\]
The map is finite and birational.  It resolves the nonnormal
cuspidal local ring without changing the function field.

\par\medskip
\noindent\textbf{Worked example 2 (an integral node with two branches).}
Assume $\operatorname{char}k\ne2$ and put
\[
 A=k[x,y]/\bigl(y^2-x^2(x+1)\bigr).
\]
The polynomial $x+1$ is not a square in $k(x)$, so the defining
polynomial is irreducible and $A$ is a domain.
Inside $K=\operatorname{Frac}(A)$ set $t=y/x$.  The defining equation
gives $t^2=x+1$, whence
\[
 x=t^2-1,\qquad y=t(t^2-1),\qquad
 A\cong k[t^2-1,t(t^2-1)]\subset k[t].
\]
Now $t$ satisfies the monic equation $T^2-(x+1)=0$ over $A$;
thus $k[t]=A+At$ is finite integral over $A$.  Because $k[t]$ is
integrally closed and shares the function field $k(t)$, it is the
normalization.  The singular point $x=y=0$ has two inverse images:
$t^2-1=0$, so $t=1$ or $t=-1$.  The two points are distinct
because $\operatorname{char}k\ne2$.  Unlike the cusp, the node has
two branches separated by normalization.

\par\medskip
\noindent\textbf{Worked example 3 (normalization of two crossing components).}
For $A=k[x,y]/(xy)$, integrality fails but $A$ is reduced.
The natural map
\[
 A\longrightarrow k[x]\times k[y],\qquad
 [f(x,y)]\longmapsto\bigl(f(x,0),f(0,y)\bigr)
\]
is injective: its kernel before quotienting is $(x)\cap(y)=(xy)$.
Its image consists exactly of pairs $(p(x),q(y))$ with
$p(0)=q(0)$.  The total ring of fractions of $A$ is
$k(x)\times k(y)$.  Put $B=k[x]\times k[y]$ in that total ring.
The idempotent $e=(1,0)$ satisfies $e^2-e=0$, so it is integral over
$A$; moreover $B=A[e]$ and is finite as an $A$-module.
Each factor of $B$ is integrally closed in its own fraction field,
so $B$ is the integral closure of $A$ in the total fraction ring.
The normalization replaces the crossing point by two origins, one
on each component.  It does not turn the reducible scheme into one
integral component.

\par\medskip
\noindent\textbf{Worked example 4 (an arithmetic normalization and its index).}
Inside $K=\mathbb Q(\sqrt5)$ consider the order
$A=\mathbb Z[\sqrt5]$.  The element
\[
 \theta=\frac{1+\sqrt5}{2}
\]
belongs to $K$ but not to $A$.  Direct calculation gives
$\theta^2-\theta-1=0$, a monic equation over $\mathbb Z$;
so $\theta$ is integral.  The ring
$B=\mathbb Z[\theta]$ contains $A$ because $\sqrt5=2\theta-1$.
For this quadratic field the standard quadratic-integer criterion
identifies $B$ as the full ring of integers $\mathcal O_K$.
As a check, the basis $1,\sqrt5$ for $A$ has discriminant $20$,
while $1,\theta$ for $B$ has discriminant $5$.
The index-discriminant formula gives
\[
 [B:A]^2=\frac{20}{5}=4,
 \qquad [B:A]=2.
\]
Thus $\operatorname{Spec}B\to\operatorname{Spec}A$ is a finite
birational normalization map, although the characteristic-zero
function field is unchanged.  This is the number-theoretic analogue
of adjoining an integral parameter to a singular affine curve.

\par\medskip
\noindent\textbf{Worked example 5 (normalization is an isomorphism off the cusp).}
Return to $A=k[t^2,t^3]$, with $x=t^2$, $y=t^3$ and normalization
$B=k[t]$.  Localize at $x$, so $t^2$ becomes invertible.
In $A_x$ we have
\[
 t=\frac{y}{x},\qquad
 t^{-1}=\frac{y}{x^2},\qquad
 A_x=k[t,t^{-1}].
\]
(The second identity follows from $y/x^2=t^3/t^4=t^{-1}$.)
Similarly $B_x=k[t,t^{-1}]$, so the localized normalization map
$A_x\to B_x$ is an isomorphism.  The closed complement of $D(x)$
on the cusp is the unique point $(x,y)$, because a prime containing
$x$ must also contain $y$ by $y^2=x^3$.
Therefore this example proves directly that normalization changes
only the cuspidal point, not the punctured curve.  More generally,
integral closure commutes with localization, so an already normal
open subset is unaffected by normalization.
